\chapter{How to use this manual}

\section{What you will build}
\begin{center}
\begin{tabularx}{\textwidth}{lXX}
\toprule
 & \textbf{Project 1: \texttt{sexconflict/}} & \textbf{Project 2: \texttt{manas\_sa/}}\\
\midrule
Question & Do the patterns of classical sex-differential selection models emerge from populations of individuals that never see the equations? & Can a mechanistic simulator explain and predict Manas Geeta Arun's sexual-antagonism results, and reconcile them with conflicting earlier studies?\\
Sources & IISER poster, Parsons 1961, Owen 1953, Havird 2019; Kidwell, Rice, Fry, Kimura & BMC Ecol Evol 2022, thesis 2022, Am Nat 2025, bioRxiv 2025 (HRI), bioRxiv 2023 (modifiers); contrast set\\
Engines & Deterministic oracle + zero-leakage \IBM\ (numba CPU, OpenCL GPU) & Analytic recursions + finite-sites multilocus \IBM\ with X/Y + hemiclone-assay emulator + GP emulator\\
Size & $\sim$5k lines of Python, 36 tests & $\sim$4k lines of Python, 8 tests, SLiM 5.2 scripts\\
Compute & minutes (quick) to $\sim$1 h (full) & $\sim$30 min (quick) to $\sim$8 h (default) on a 4-core laptop\\
\bottomrule
\end{tabularx}
\end{center}

\section{Reading paths}
\begin{itemize}
\item \textbf{For the big picture first:} read Chapter~\ref{ch:summary} (plain language, no equations), then Chapter~\ref{ch:examples} (every model and algorithm worked with small numbers).
\item \textbf{To rebuild everything:} read the chapters in order. Chapter~\ref{ch:background} gives the theory you must understand before coding. Chapters~\ref{ch:build1} and~\ref{ch:build2} are the build instructions. Keep Chapter~\ref{ch:trouble} open while you work.
\item \textbf{To use the results only:} read Chapter~\ref{ch:sources} (what the papers claim), Chapter~\ref{ch:results} (what the models say) and Chapter~\ref{ch:heur} (what cannot be concluded).
\item \textbf{To extend the simulator:} read Chapter~\ref{ch:design} (the rules that keep the work honest) and Chapter~\ref{ch:heur} (nuances that bit us).
\end{itemize}

\section{Conventions}
\begin{keybox}[Claim types and verdict words]
\begin{itemize}
\item \claimA\ means the value is reported in the paper.
\item \claimB\ means we reproduced it, for example by recomputing from raw data or re-running the authors' code.
\item \claimC\ means it is our extrapolation or a model prediction.
\end{itemize}
Comparisons between model and paper use only these verdicts: \emph{matches sign}, \emph{matches ordering}, \emph{misses magnitude}, \emph{not identified}. Designed null checks use \emph{null holds} or \emph{null fails}.

Never write ``the model confirms'' or ``validates''. A simulator can only be \emph{consistent with}, or \emph{inconsistent with}, a result under stated assumptions.
\end{keybox}
\begin{itemize}
\item Code listings are taken from the working repository, sometimes shortened. File paths are relative to the project folder.
\item Notation:
  \begin{itemize}
  \item $\rmf$ is the intersexual genetic correlation for fitness computed on genotypes, as in the HRI preprint.
  \item $\rwg$ is the same quantity computed from hemiclone line means (BMC 2022).
  \item $\lambda$ is the population rescaling factor.
  \item MB / E / FB are the male-biased, equal and female-biased adult sex ratios.
  \end{itemize}
\end{itemize}

\section{Prerequisites}
\begin{itemize}
\item \textbf{Knowledge.}
  \begin{itemize}
  \item Population genetics at the level of Hartl \& Clark: selection recursions, drift, linkage disequilibrium and recombination.
  \item Basic quantitative genetics: breeding values, genetic covariance, $h^2$.
  \item Python with numpy.
  \end{itemize}
  The specialised material (hemiclonal analysis, Hill--Robertson interference, the Robertson/FTNS decomposition, rescaling, numba bit tricks) is explained here.
\item \textbf{Software.} Python 3.14 with the pinned packages of Appendix~\ref{app:env}. Optionally:
  \begin{itemize}
  \item pyopencl for AMD/other GPUs;
  \item SLiM 5.2 for cross-checks;
  \item MiKTeX or TeX Live for this document.
  \end{itemize}
\item \textbf{Hardware.} Anything with 4+ cores and 8\,GB RAM. The reference machine was a Ryzen 5 2500U (4C/8T), 32\,GB, with a Radeon RX 560X, shared with other workloads.
\end{itemize}

\section{The shortest path to a working system}
\begin{lstlisting}[style=sh]
python -m venv .venv
.venv/Scripts/pip install -r requirements.txt            # project 1
.venv/Scripts/pip install -r manas_sa/requirements.txt   # project 2
.venv/Scripts/python -m pytest -q tests manas_sa/tests
.venv/Scripts/python scripts/run_all.py --profile quick               # project 1, minutes
cd manas_sa && ../.venv/Scripts/python scripts/run_all.py --profile quick  # project 2, ~30 min
../.venv/Scripts/python scripts/make_comparison.py        # comparison table from cached results
\end{lstlisting}
On Linux or macOS, replace \code{.venv/Scripts/} with \code{.venv/bin/}.
